\subsection{各种自由对象之间的关系}

由于自由幺半群 \(\mathrm{Mo}(\mathbf{X})\) 是一个广群，第 1 号中的命题 1 表明存在一个同态 \(\lambda: \mathrm{M}(\mathbf{X}) \to \mathrm{Mo}(\mathbf{X})\) ，它在 \(\mathbf{X}\) 上的限制是恒等映射。类似地，由于自由群 \(\mathrm{F}(\mathbf{X})\) 是幺半群， \(\mathbf{X}\) 的恒等映射可扩张为同态 \(\mu: \mathrm{Mo}(\mathbf{X}) \to \mathrm{F}(\mathbf{X})\) （第 2 号，命题 3）。由第 4 号命题 6 及第 5 号命题 7， \(\mu\) 是单射的。类似地，第 7 号的命题 10 和第 5 号的命题 8 表明存在同态 \(\nu: \mathrm{Mo}(\mathbf{X}) \to \mathbf{N}^{(\mathbf{X})}\) 和 \(\pi: \mathrm{F}(\mathbf{X}) \to \mathbf{Z}^{(\mathbf{X})}\) ，其特征是 \(\nu(x) = \delta_x\) 和 \(\pi(x) = \delta_x\) 对所有 \(x \in \mathbf{X}\) 成立。若 \(\iota\) 是 \(\mathbf{N}^{(\mathbf{X})}\) 到 \(\mathbf{Z}^{(\mathbf{X})}\) 中的单射，则这两个同态 \(\iota \circ \nu\) 和 \(\pi \circ \mu\) （都把 \(\mathrm{Mo}(\mathbf{X})\) 映入 \(\mathbf{Z}^{(\mathbf{X})}\) ）在 \(\mathbf{X}\) 上重合，因此 \(\iota \circ \nu = \pi \circ \mu\) 。这种情况可以用下面的交换图来概括：

\[
\begin{array}{c} \mathrm{M} (\mathrm{X}) \xrightarrow {\lambda} \mathrm{Mo} (\mathrm{X}) \xrightarrow {\nu} \mathbf {N} ^ {(\mathrm{X})} \\ \mu \Biggl \downarrow \\ \mathrm{F} (\mathrm{X}) \xrightarrow {\pi} \mathbf {Z} ^ {(\mathrm{X})}. \end{array}
\]

同态 \(\lambda\) , \(\mu\) , \(\nu\) 和 \(\pi\) 将被称为典范的。

设 \(w\) 的序列 \(\mathbf{M}(\mathbf{X})\) ；通过对 \(l(w)\) 进行归纳可以立即证明，单词 \(\lambda(w)\) 的长度等于 \(w\) 的长度。此外，

\[
\nu (x _ {1} \dots x _ {n}) = \sum_ {i = 1} ^ {n} \delta_ {x _ {i}}\tag{21}
\]

对于 \(x_1, \ldots, x_n\) 在X中，因此根据(13)和(14)， \(|\nu(x_1 \ldots x_n)| = n\) 。换句话说，

\[
\left| \nu (u) \right| = l (u) \quad (u \in \operatorname{Mo} (\mathbf {X})).\tag{22}
\]

命题11. 典范同态 \(\nu\) 于 \(\operatorname{Mo}(\mathbf{X})\) 到 \(\mathbf{N}^{(\mathbf{X})}\) 是满射。设 \(w = x_1 \ldots x_n\) 和 \(w' = x_1' \ldots x_m'\) 是 \(\operatorname{Mo}(\mathbf{X})\) 的两个元素；为了使 \(\nu(w) = \nu(w')\) ，必要且充分的条件是 \(m = n\) 并且存在一个置换 \(\sigma \in \mathfrak{S}_n\) 使得 \(x_i' = x_{\sigma(i)}\) 对于 \(1 \leqslant i \leqslant n\) .

\(\nu\) 的像是 \(\mathbf{N}^{(\mathbf{x})}\) 的一个子幺半群 I，它包含元素 \(\delta_x\) （对 \(x \in \mathbf{X}\) ）。公式 (12)（第 7 号）表明， \(\mathbf{N}^{(\mathbf{x})}\) 由族 \((\delta_x)_{x \in \mathbf{X}}\) 生成，其中 \(\mathbf{I} = \mathbf{N}^{(\mathbf{x})}\) 。因此 \(\nu\) 是满射。

若 \(m = n\) 且 \(x_{i}^{\prime} = x_{\sigma (i)}\) （对 \(\mathrm{l}\leqslant i\leqslant n\) ），则

\[
\nu (w ^ {\prime}) = \sum_ {i = 1} ^ {n} \delta_ {x _ {i} ^ {\prime}} = \sum_ {i = 1} ^ {n} \delta_ {x _ {\sigma (i)}} = \sum_ {i = 1} ^ {n} \delta_ {x _ {i}} = \nu (w)
\]

由公式 (21) 和交换性定理（§ 1，第 5 号，定理 2）得出。

反之，假设 \(\nu(w)\) 和 \(\nu(w')\) 等于同一个元素 \(\alpha\) （属于 \(\mathbf{N}^{(\mathbf{x})}\) ）；由公式 (22)， \(n = |\alpha| = m\) 。对所有 \(x \in \mathbf{X}\) ，设 \(\mathrm{I}_x\) （相应地 \(\mathrm{I}'_x\) ）为整数 \(i\) 组成的集合，满足 \(1 \leqslant i \leqslant n\) 且 \(x_i = x\) （相应地 \(x_i' = x\) ）。因此 \((\mathrm{I}_x)_{x \in \mathbf{X}}\) 和 \((\mathrm{I}'_x)_{x \in \mathbf{X}}\) 是区间 \([1, n]\) （\(\mathbf{N}\) 中的）的分割；进一步，公式 \(\alpha = \sum_{i=1}^{n} \delta_{x_i}\) 表明 \(\alpha(x)\) 是 \(I_x\) 的基数；类似地，公式 \(\alpha = \sum_{i=1}^{\infty} \delta_{x'_i}\) 表明 \(\alpha(x)\) 是 \(\mathbf{I}_x'\) 的基数。因此存在一个置换 \(\sigma\) （\([1, n]\) 的），使得 \(\sigma(\mathbf{I}_x') = \mathbf{I}_x\) 对所有 \(x \in \mathbf{X}\) 成立，即 \(x'_i = x_{\sigma(i)}\) 对 \(i = 1, \ldots, n\) 成立。

注记。设 S 是 X 的一个子集。回顾我们已将 M(S) 与 M(X) 的一个子广群等同，将 Mo(S) 与 Mo(X) 的一个子幺半群等同，并且将 \(\mathbf{N}^{(S)}\) 与 \(\mathbf{N}^{(X)}\) 的一个子幺半群等同。则

\[
\mathbf {M} (\mathrm{S}) = \lambda^ {- 1} (\mathbf {M o} (\mathrm{S})).\tag{23}
\]

显然 \(\lambda (\mathbf{M}(\mathbf{S}))\subset \mathbf{Mo}(\mathbf{S})\) 。设 \(w\in \lambda^{-1}(\mathbf{Mo}(\mathbf{S}))\) ；我们通过对 \(l(w)\) 作归纳来证明 \(w\in \mathbf{M}(\mathbf{S})\) 。当 \(l(w) = 1\) 时这是显然的。若 \(l(w) > 1\) ，我们可以写成 \(w = w_{1}w_{2}\) ，其中 \(w_{1},w_{2}\in \mathbf{M}(\mathbf{X}), l(w_{1}) <   l(w),l(w_{2}) <   l(w)\) 。则 \(\lambda (w_1)\lambda (w_2)\in \mathbf{Mo}(\mathbf{S})\) ，因此 \(\lambda (w_1)\in \mathbf{Mo}(\mathbf{S})\) 且 \(\lambda (w_2)\in \mathbf{Mo}(\mathbf{S})\) ，从而由归纳假设 \(w_{1}\in \mathbf{M}(\mathbf{S})\) 且 \(w_{2}\in \mathbf{M}(\mathbf{S})\) ，最后 \(w\in \mathbf{M}(\mathbf{S})\) 。

此外

\[
\mathbf {M o} (\mathrm{S}) = \nu^ {- 1} (\mathbf {N} ^ {(\mathrm{S})}).\tag{24}
\]

这由公式（21）直接得出。

进一步， \(\mathbf{N}^{(\mathbf{S})}\) 是 \(\mathbf{N}^{(\mathbf{X})}\) 中支集包含于 S 的元素组成的集合；若 \((\mathrm{S}_i)_{i\in I}\) 是 X 的一个子集族，其交为 S，则 \(\mathbf{N}^{(\mathbf{S})} = \bigcap_{i\in I}\mathbf{N}^{(\mathbf{S}_i)}\) 且公式 (23) 和 (24) 蕴含

\[
\mathbf {M} (\mathrm{S}) = \bigcap_ {i \in \mathrm{I}} \mathbf {M} (\mathrm{S} _ {i}), \quad \mathbf {M o} (\mathrm{S}) = \bigcap_ {i \in \mathrm{I}} \mathbf {M o} (\mathrm{S} _ {i}).\tag{25}
\]

